%% Abridged CV (3-page version) - Prof. Muaaz Bhamjee
%% -------------------------------------------------------------------
%% Abridged CV for funding applications that impose a 3-page limit
%% (e.g. Oppenheimer Memorial Trust New Frontiers Research Award).
%% Reuses the styling of cv_bhamjee_main.tex, but with CURATED content
%% and no auto-injected publication list, so that it fits in 3 pages.
%% Built as a plain 'latex' document via build.py (see DOCUMENTS registry).
%% -------------------------------------------------------------------

\documentclass[11pt, a4paper]{article}

% ---- Packages (as per cv_bhamjee_main.tex) ----
\usepackage[a4paper, top=1.5cm, bottom=1.5cm, left=2.0cm, right=2.0cm]{geometry}
\usepackage{hyperref}
\usepackage{enumitem}
\usepackage{titlesec}
\usepackage{parskip}
\usepackage{array}
\usepackage{booktabs}
\usepackage{longtable}
\usepackage[T1]{fontenc}
\usepackage[utf8]{inputenc}
\usepackage{textcomp}
\usepackage{xcolor}
\usepackage{amsmath}

\hypersetup{
    colorlinks=true,
    linkcolor=blue!70!black,
    urlcolor=blue!70!black,
    citecolor=blue!70!black
}

% ---- Section formatting (slightly tighter spacing to hold 3 pages) ----
\titleformat{\section}
  {\large\bfseries\uppercase}
  {}
  {0em}
  {}
  [\titlerule]
\titlespacing{\section}{0pt}{7pt}{3pt}

\setlist[itemize]{noitemsep, topsep=1pt, leftmargin=1.4em}
\setlength{\parindent}{0pt}
\setlength{\parskip}{2.5pt}

% ===============================================================
\begin{document}

% ---- Header ----
\begin{center}
  {\LARGE\bfseries Muaaz Bhamjee}\\[3pt]
  {\large Associate Professor, Department of Mechanical and Aeronautical Engineering, University of Pretoria}\\[4pt]
  Mobile: +27 76 476 4434 \quad Office: +27 12 420 5366 \quad Email: \href{mailto:muaaz.bhamjee@up.ac.za}{muaaz.bhamjee@up.ac.za}\\[2pt]
  Language: English \quad Gender: Male \quad Citizenship: South African\\[2pt]
  Google Scholar h-index: 68 \quad Scopus h-index: 35 \quad Web of Science h-index: 31 \quad ORCID: \href{https://orcid.org/0000-0002-2697-4589}{0000-0002-2697-4589}
\end{center}

% ===============================================================
\section{Profile}

I am Professor Muaaz Bhamjee, Associate Professor in the Department of Mechanical and Aeronautical Engineering at the University of Pretoria. 
My research programme spans computational fluid dynamics and granular mechanics, machine learning and AI applied to fluid dynamics, geospatial and climate systems, 
experimental high-energy physics through the ATLAS Collaboration at CERN, and emerging quantum computing and quantum sensing programmes that I am establishing at UP. 
I hold a BIng and MIng (Mechanical Engineering), a BSc Honours (Applied Mathematics), and a DIng (Computational Fluid and Granular Dynamics), all from the University of Johannesburg, 
where I served as a Lecturer and Senior Lecturer before transitioning to IBM Research Africa as a Staff Research Scientist and subsequently to my current position at UP.

The unifying thread across all of these programmes is \textbf{physics-informed computational simulation of transport phenomena}, the development and validation of numerical models 
that resolve how physical systems transport energy, matter, and information, applied wherever the physics demands it: from granular particles in a 
hydrocyclone to acoustic waves in a diseased lung to Higgs bosons decaying in the ATLAS detector.

I established CERG-FLUX Lab (Fluids, Learning, and Uncertainty in compleX systems) at the University of Pretoria, 
to host my research programme across CFD, Lattice Boltzmann Methods, Scientific Machine Learning, ATLAS/CERN particle physics, and quantum technologies. 
The group currently comprises three MEng students, eight PhD students, and postdoctoral fellows in pipeline. CERG-FLUX Lab is involved in quantum computing and quantum sensing, 
with applications to fluid dynamics, particle physics, and mineral separation - housed in the UP Quantum Science and Technology (UPQuST) initiative, which I co-founded in 2026 
and I am a Principal Investigator. UPQuST is the 6th node of the South African Quantum Technology Initiative (SA QuTI), a national network of quantum research groups.

% ===============================================================
\section{Education}
\textbf{Doctor of Engineering (DIng, Mechanical Engineering Science)} \hfill 2011--2017\\
\textit{University of Johannesburg}\\
Thesis: Modelling of the Multiphase Interactions in a Hydrocyclone using Navier-Stokes and Lattice Boltzmann based Computational Approaches.

\smallskip
\textbf{Master of Engineering (MIng, Mechanical Engineering Science), \textit{Cum Laude}} \hfill 2009--2011\\
\textit{University of Johannesburg}\\
Dissertation: A Computational Fluid Dynamics and Experimental Investigation of an Airflow Window.

\smallskip
\textbf{BSc Honours (Applied Mathematics), \textit{Cum Laude}} \hfill 2008--2010\\
\textit{University of Johannesburg}

\smallskip
\textbf{Bachelor of Engineering (BIng, Mechanical Engineering Science)} \hfill 2004--2008\\
\textit{University of Johannesburg}

% ===============================================================
\section{Academic and Professional Positions}
\textbf{Associate Professor}, \textit{University of Pretoria} \hfill Feb 2023 -- Present\\
Founder and Principal Investigator, CERG-FLUX Lab. Teaching (Computational Fluid Dynamics (MKM411) and Impact of Engineering Activity and Group Work (MIA 320)), postgraduate supervision, and research in computational mechanics, scientific machine learning, and quantum technologies. UP ATLAS / CERN Institutional Representative and Team Leader.

\smallskip
\textbf{Staff Research Scientist}, \textit{IBM Research -- Africa} \hfill Feb 2023 -- Dec 2024\\
Data-driven and machine-learning approaches to climate; challenge / project lead; patent ideation; student and intern mentorship.

\smallskip
\textbf{Senior Research Associate (Visiting)}, \textit{University of Johannesburg} \hfill Jul 2023 -- Dec 2024\\
Research and postgraduate supervision.

\smallskip
\textbf{Senior Lecturer}, \textit{University of Johannesburg} \hfill May 2017 -- Jan 2023\\
Course co-ordination and lecturing; postgraduate supervision; CFD and engineering-education research; former co-ordinator of the MEng (Sustainable Energy) programme. SA-CERN ATLAS Principal Investigator.

\smallskip
\textbf{Lecturer}, \textit{University of Johannesburg} \hfill Jan 2015 -- Apr 2017\\
Course co-ordination and lecturing; final-year project supervision; co-supervision of Masters students; CFD and engineering-education research.

\smallskip
\textbf{Junior Engineer (CFD Analyst)}, \textit{Hatch Africa (Pty) Ltd} \hfill 2008 -- 2009\\
CFD modelling for the metals and mineral-processing environment; user-defined functions; technical reporting.

% ===============================================================
\section{Research Leadership and Professional Affiliations}
\begin{itemize}
  \item Founder and Principal Investigator, CERG-FLUX Lab, University of Pretoria (2026--present).
  \item Principal Investigator, SA-CERN ATLAS Project; University of Pretoria ATLAS / CERN Institutional Representative and Team Leader.
  \item Registered Professional Engineer (PrEng), Engineering Council of South Africa (ECSA).
  \item Vice-President, SAAM Executive Committee; Vice-President, South African National IUTAM Committee.
  \item Former: Secretary, SAAM Executive Committee; Secretary, South African National IUTAM Committee.
  \item Member, South African Institute of Mechanical Engineering (SAIMechE) and the South African Association for Theoretical and Applied Mechanics (SAAM).
\end{itemize}

% ===============================================================
\section{Noteworthy Research Grants and Funding}
\begin{itemize}
  \item \textbf{2026 -- UPQuST / SA QuTI-linked grants:} Hybrid Quantum-Classical Accelerated CFD (R180,000); Quantum Machine Learning for Rare Particle Searches in HEP (R180,000); Quantum-Magnetic Sensors for Particle Detection in Mineral Separation (R400,000).
  \item \textbf{2022 -- UJ URC Research Equipment Grant:} Preventing COVID-19 with Mathematical Modelling, Digital Twins, AI and 4IR Engineered Solutions (R284,517).
  \item \textbf{2021 -- UJ GES COVID-19 Grant:} same programme (R80,000).
  \item \textbf{2020--2021 -- European Synchrotron Radiation Facility (ESRF):} Equipment for the BEATS beamline at SESAME and ESRF beamlines (\texteuro8,750).
  \item \textbf{2017--2019 -- NRF Thuthuka Funding Instrument (Post-PhD Track):} Modelling of Multiphase Flow in Process Equipment (TTK160615171464; R602,000).
  \item \textbf{Patent:} Dynamically Forecasting High-Resolution Air Temperature in Real-Time Using Multiple Sources -- ZA2023/09448 (granted 30 July 2024); USA pending (US18/200,050).
  \item \textbf{CHPC high-performance computing allocations:} Multiphase Flow Modelling and Deep Learning (MECH1739, active); Modelling of Multiphase Flow in Process Equipment (MECH0837).
\end{itemize}

% ===============================================================
\section{Research Output and Postgraduate Supervision}
Publication record indexed on Google Scholar (h-index 68), Scopus (h-index 35) and Web of Science (h-index 31), spanning journal articles, ATLAS Collaboration papers, conference papers, a book chapter and a granted patent. Full publication list available separately. Postgraduate supervision: 9 Masters graduated; currently main or co-supervisor of 8 doctoral and 3 Masters candidates.

% ===============================================================
\section{Current Doctoral Supervision}
\begin{itemize}
  \item M. Chambakata, ``Lattice Boltzmann Method for Modelling Two-Phase Flow'' (main supervisor).
  \item A.Y. Akinbowale, ``Quantum Machine Learning for Rare Particle Signatures in High-Dimensional Physics Data'' (main supervisor).
  \item M. Ramokali, ``Hybrid Classical-Quantum Computing and Machine Learning Methods for Computational Fluid Dynamics'' (main supervisor).
  \item M.T. Saphira Agnes, ``Fluid-Structure Interactions in Pathological Arteries: Numerical and Experimental Models'' (main supervisor).
  \item M.S.W. Potgieter, ``Hybrid Lattice Boltzmann -- Machine Learning Methods for Two-Phase Flow'' (main supervisor).
  \item A. Essop, ``A Novel Numerical Scheme for High-Speed Flows'' (main supervisor).
  \item T. Jooma Abbajee, ``Graph Dynamical Reservoir Computing'' (co-supervisor).
  \item M.P. Connell, ``Search for Higgs Boson Decays via Dark Bosons to a Four-Lepton Final State'' (co-supervisor).
\end{itemize}

% ===============================================================
\section{Selected Postgraduate Supervision (Graduated)}
\begin{itemize}
  \item T. Jooma Abbajee, ``Reservoir Computing: A Tool for Predicting Chaotic Dynamical Systems,'' University of Johannesburg, 2026 (co-supervisor).
  \item K.M. Makhanya, ``The Use of CFD for Assessing Flow-Induced Acoustics to Diagnose Lung Conditions,'' University of Johannesburg, 2025 (co-supervisor).
  \item T.F. Mokoena, ``Experimental Hutch Equipment for the BEATS Beamline at SESAME,'' University of Johannesburg, 2022 (main supervisor).
  \item T. Ramokoka, ``Fluid and Particulate Flow Through a Ventriculoperitoneal Shunt in a Variable-Temperature Environment,'' University of Johannesburg, 2022 (sole supervisor).
  \item A. Potgieter, ``Modelling of a Heated Gas-Solid Fluidised Bed using Eulerian Based Models,'' University of Johannesburg, 2020, \textit{Cum Laude} (main supervisor).
\end{itemize}

% ===============================================================
\section{Selected Publications}
\begin{itemize}
  \item ``Assessment and characterization of the impact of pulmonary pathology on flow-induced acoustics using computational fluid dynamics,'' \textit{Biomedical Signal Processing and Control}, \textbf{120}, 110018, (2026).
  \item ``Characterisation of the Thermoflow due to the Dry Nitrogen Flushing Scheme in the ATLAS Inner Tracker using Computational Fluid Dynamics,'' \textit{arXiv preprint}, \textbf{arXiv:2605.08331} (2026). (under review, \textit{Nucl.\ Instrum.\ Methods Phys.\ Res.\ A})
  \item ATLAS Collaboration, ``Characterising the Higgs boson with ATLAS data from the LHC Run-2,'' \textit{Physics Reports} \textbf{1116}, 4--56 (2025).
  \item ``Reservoir Computing for Predicting Chaotic Dynamical Systems,'' \textit{69th Annual Conference of the South African Institute of Physics (SAIP 2025)}, 1026--1031 (2025).
  \item ``Infectiousness model of expelled droplets exposed to ultraviolet germicidal irradiation,'' \textit{Computers \& Fluids} \textbf{275}, 106242 (2024).
  \item ATLAS Collaboration, ``Observation of quantum entanglement with top quarks at the ATLAS detector,'' \textit{Nature} \textbf{633}, 542--547 (2024).
  \item ``Detection and Characterization of Urban Heat Islands with Machine Learning,'' \textit{IGARSS 2024 -- IEEE International Geoscience and Remote Sensing Symposium}, Athens, 1693--1699 (2024).
  \item ``The Modification of the Dynamic Behaviour of Cyclonic Flow in a Hydrocyclone under Surging Conditions,'' \textit{Mathematical and Computational Applications} \textbf{27}(6), 88 (2022).
  \item ``Experimental and CFD Investigation of a Hybrid Solar Air Heater,'' \textit{Solar Energy} \textbf{195}, 413--428 (2020).
  \item ``An Experimentally Validated Mathematical and CFD Model of a Supply Air Window: Forced and Natural Flow,'' \textit{Energy and Buildings} \textbf{57}, 289--301 (2013).
\end{itemize}

% ===============================================================
\section{Selected Professional Service}
\begin{itemize}
  \item Journal reviewing: \textit{Physics of Fluids} (AIP), \textit{Applied Energy} (Elsevier), \textit{Energy and Buildings} (Elsevier), \textit{Environmental Science and Pollution Research} (Springer), and the \textit{SAIMechE R\&D Journal}; conference reviewing for SAIP and SACAM, and ICLR 2024 workshop proposals.
  \item Conference organising committees: Kruger2026 (7th Biennial Workshop on Discovery Physics at the LHC, SA-CERN); SAIMechE Central Branch Conference 2025 (Wits); SMPM 2019; SACAM 2012.
  \item Postgraduate examination: external examiner for doctoral and Masters candidates at the Universities of Cape Town, Johannesburg and KwaZulu-Natal; internal examiner at the University of Pretoria.
\end{itemize}

\end{document}
